\section{Simulationsstudien}

Zur Überprüfung des chemischen Analogcomputermodells wurden zwei
Simulationsstudien durchgeführt. In der ersten Studie wird untersucht, ob
das Modell für verschiedene Eingabewerte \(s\) die korrekten stationären
Ausgabewerte berechnet. In der zweiten Studie wird bestimmt, nach welcher
Modellzeit alle relevanten Ausgabespezies eine Genauigkeit von mindestens
\(99.9\,\%\) erreicht haben.

Alle Simulationen wurden als deterministische Zeitverläufe in COPASI
durchgeführt. Die Eingabe wurde jeweils durch die Konzentration der
Spezies \(S\) kodiert. Die Spezies \(Q\), \(F\) und \(\Theta_n\) wurden mit
Anfangskonzentration \(0\) initialisiert.

\subsection{Studie 1: Validierung der Berechnung für verschiedene Eingaben}

In der ersten Simulationsstudie wurde überprüft, ob das Modell für
verschiedene Werte der Eingabespezies \(S\) die erwarteten stationären
Werte liefert. Dazu wurde ein Parameter-Scan über die Anfangskonzentration
von \(S\) durchgeführt. Für jeden Eingabewert wurde ein Zeitverlauf bis zur
Modellzeit $T = \text{30}$ mit einer Intervallgröße von \(\Delta t = \text{0.1}\) Zeiteinheiten und 300 Intervallen simuliert. Am Ende der Simulation wurden die Werte der Spezies \(Q\), \(F\) und \(\Theta_n\) mit den theoretisch erwarteten Gleichgewichtswerten verglichen. Die Ergebnisse des Parameter-Scans sind in Tabelle in\ref{tab:studie1_parameter_scan} dargestellt. (TODO: adde theoretische Werte in die Tabelle)

\begin{table}[h]
    \centering
    \begin{tabular}{c|c|c|c|c}
        \(s=[S]_0\) &
        \(t\) &
        \([Q](t)\) &
        \([F](t)\) &
        \([\Theta_n](t)\) \\
        \hline
        0.5 & 30 & 0.25 & 0.429112 & 0.819546 \\
        1	&30	&1&	1.71645	& 0.819546 \\
        2	&30	&4	&6.86578	&0.819546 \\
        3	&30	&9	&15.448	&0.819546 \\
        5	&30	&25	&42.9112	&0.819546 \\
    \end{tabular}
    \caption{Endwerte des Parameter-Scans zur Modellzeit \(t=30\).}
    \label{tab:studie1_parameter_scan}
\end{table}

Die Simulationsergebnisse zeigen, dass die Spezies \(Q\) für alle
betrachteten Eingaben gegen \(s^2\) konvergiert. Die Ausgabespezies \(F\)
nähert sich entsprechend dem theoretischen Wert \(A_{\max}(s)\). Die
normierte Winkelausgabe \(\Theta_n\) bleibt unabhängig von \(s\) konstant,
wie es aufgrund der mathematischen Herleitung zu erwarten ist. Damit
bestätigt die erste Studie die korrekte Implementierung der
Berechnungsvorschriften.

\subsection{Studie 2: Bestimmung der chemischen Rechenzeit}

In der zweiten Simulationsstudie wurde untersucht, nach welcher Modellzeit
alle relevanten Ausgabespezies eine Genauigkeit von \(99.9\,\%\) erreichen.
Dazu wurde ein repräsentativer Eingabewert $s = 2$ verwendet und ein Zeitverlauf mit feiner Schrittweite simuliert.

Ein Ergebnis gilt als mit \(99.9\,\%\) Genauigkeit vorliegend, wenn der
relative Fehler höchstens \(0.001\) beträgt. Für eine Ausgabespezies \(X\)
mit Zielwert \(X^*\) bedeutet dies

\begin{equation}
    \frac{|X^* - X(t)|}{X^*}
    \leq
    0.001.
\end{equation}

Da alle betrachteten Spezies monoton von \(0\) gegen ihren Zielwert
konvergieren, kann äquivalent geprüft werden, ob

\begin{equation}
    X(t)\geq 0.999X^*
\end{equation}

gilt.

Die verwendeten Zielwerte und die entsprechenden \(99.9\,\%\)-Schwellen
sind in Tabelle \ref{tab:studie2_schwellen} dargestellt.

\begin{table}[h]
    \centering
    \begin{tabular}{c|c|c}
        Spezies & Zielwert \(X^*\) & \(99.9\,\%\)-Schwelle \(0.999X^*\) \\
        \hline
        \(Q\) & 4 & 3.996 \\
        \(F\) & 6.865784432 & 6.858918648 \\
        \(\Theta_n\) & 0.819546 & 0.8187269 \\
    \end{tabular}
    \caption{Zielwerte und \(99.9\,\%\)-Schwellen für die Bestimmung der
    chemischen Rechenzeit.}
    \label{tab:studie2_schwellen}
\end{table}

Anschließend wurde aus dem simulierten Zeitverlauf jeweils die erste
Modellzeit bestimmt, zu der die entsprechende Spezies ihre
\(99.9\,\%\)-Schwelle erreicht. Die Ergebnisse sind in Tabelle
\ref{tab:studie2_rechenzeit} zusammengefasst.

\begin{table}[h]
    \centering
    \begin{tabular}{c|c|c}
        Spezies & erste Modellzeit mit \(X(t)\geq0.999X^*\) & Bemerkung \\
        \hline
        \(Q\) & 7 & Quadratwert \(s^2\) \\
        \(F\) & 9.3 & maximaler Flächeninhalt \\
        \(\Theta_n\) & 7 & normierter optimaler Winkel \\
    \end{tabular}
    \caption{Bestimmung der Modellzeit bis zum Erreichen einer
    Genauigkeit von \(99.9\,\%\).}
    \label{tab:studie2_rechenzeit}
\end{table}

Die chemische Rechenzeit des Gesamtmodells wird durch die langsamste der
betrachteten Ausgabespezies bestimmt. Damit gilt

\begin{equation}
    t_{\mathrm{chem}}
    =
    \max
    \left\{
        t_Q,\,
        t_F,\,
        t_{\Theta_n}
    \right\}.
\end{equation}

Aus den Simulationsergebnissen folgt

\begin{equation}
    t_{\mathrm{chem}}
    =
    \text{9.3}.
\end{equation}

Die langsamste Konvergenz tritt bei der Flächenausgabe \(F\) auf. Dies ist
dadurch erklärbar, dass \(F\) nicht direkt aus der Eingabe \(S\), sondern
über die Zwischenspezies \(Q\) berechnet wird. Das Modell besitzt also für
die Flächenberechnung eine Kaskadenstruktur

\begin{equation}
    S \longrightarrow Q \longrightarrow F.
\end{equation}

Dadurch reagiert \(F\) verzögert auf die Annäherung von \(Q\) an seinen
Gleichgewichtswert.

Ergänzend kann die Rechenzeit analytisch aus der Struktur des
Differentialgleichungssystems nachvollzogen werden. Für \(Q\) und
\(\Theta_n\) ergibt sich jeweils eine Konvergenz der Form

\begin{equation}
    X(t)=X^*\left(1-e^{-t}\right).
\end{equation}

Die Bedingung \(X(t)\geq0.999X^*\) führt auf

\begin{equation}
    t\geq \ln(1000)\approx6.91.
\end{equation}

Für die Flächenausgabe \(F\) ergibt sich aufgrund der Kaskade

\begin{equation}
    F(t)
    =
    F^*
    \left(
        1-(1+t)e^{-t}
    \right).
\end{equation}

Die Bedingung \(F(t)\geq0.999F^*\) führt auf

\begin{equation}
    (1+t)e^{-t}\leq0.001.
\end{equation}

Daraus ergibt sich näherungsweise

\begin{equation}
    t\approx9.24.
\end{equation}

Dieses Ergebnis stimmt mit der simulativ bestimmten chemischen Rechenzeit
überein. Somit liegen alle berechneten Ausgaben nach ungefähr
\(\text{9.3}\) Modellzeiteinheiten mit einer Genauigkeit
von mindestens \(99.9\,\%\) vor.